\begin{proof}[Proof of \cref{obj:thm:degree-frontier}]
\begin{enumerate}
\item
Apply \cref{obj:thm:bounded-outcome-degree-frontier} with the fixed \(\beta\) and \(p\).  It gives constants
\(\alpha_{\beta,p},\Gamma_{\beta,p}\in\mathbb R\), depending only on \((\beta,p)\), such that
\[
0<\alpha_{\beta,p}\le \Gamma_{\beta,p}.
\]
Set
\[
c_{\beta,p}:=\min\{1,\alpha_{\beta,p}^{2}\},
\qquad
C_{\beta,p}:=\max\{1,\Gamma_{\beta,p}^{2}\}.
\]
Then \(0<c_{\beta,p}\) and \(c_{\beta,p}\le C_{\beta,p}\).  Now fix \(n,d\in\mathbb N\) and \(B\in\mathbb R\) with \(1\le n\), \(d\le n\), and \(0\le B\).  Let \(D_{n,p}\) be the product Bernoulli law on \(\{0,1\}^{V_n}\),
\[
D_{n,p}(z)=\prod_{j\in V_n}p^{z_j}(1-p)^{1-z_j}.
\]
This is the product Bernoulli design with common probability \(p\), so \cref{obj:thm:bounded-outcome-degree-frontier} applies to \(D_{n,p}\).  Write
\[
x_{n,d}:=\frac{d\binom{d}{k_\star(d,\beta,p)}}{n},
\qquad
y_{n,d}:=\frac{dA_d}{n}.
\]
The same conclusion supplies
\[
\alpha_{\beta,p}\,d\binom{d}{k_\star(d,\beta,p)}
\le dA_d
\le
\Gamma_{\beta,p}\,d\binom{d}{k_\star(d,\beta,p)}.
\]
Since \(n\ge1\), division by \(n\) gives
\[
\alpha_{\beta,p}x_{n,d}\le y_{n,d}\le \Gamma_{\beta,p}x_{n,d}.
\]
The case \(d=0\) is included in these displayed inequalities; then \(x_{n,d}=y_{n,d}=0\). % lean: degree_frontier

\item
The elementary rescalings used to pass from \(y_{n,d}\) to \(x_{n,d}\) are
\[
c_{\beta,p}\min\{1,x_{n,d}\}
\le
\alpha_{\beta,p}\min\{1,\alpha_{\beta,p}x_{n,d}\},
\]
and
\[
\Gamma_{\beta,p}\min\{1,\Gamma_{\beta,p}x_{n,d}\}
\le
C_{\beta,p}\min\{1,x_{n,d}\}.
\]
For the first inequality, split into \(x_{n,d}\le1\) and \(x_{n,d}>1\), and inside each case split again according as \(\alpha_{\beta,p}x_{n,d}\le1\) or \(\alpha_{\beta,p}x_{n,d}>1\); use \(c_{\beta,p}\le\alpha_{\beta,p}^{2}\), \(c_{\beta,p}\le\alpha_{\beta,p}\), and \(x_{n,d}\ge0\).  The second inequality follows from the same two splits, using
\[
\Gamma_{\beta,p}\le C_{\beta,p},
\qquad
\Gamma_{\beta,p}^{2}\le C_{\beta,p},
\qquad
x_{n,d}\ge0.
\]
Together with \(\alpha_{\beta,p}x_{n,d}\le y_{n,d}\le \Gamma_{\beta,p}x_{n,d}\), these imply
\[
c_{\beta,p}\min\{1,x_{n,d}\}
\le
\alpha_{\beta,p}\min\{1,y_{n,d}\},
\]
and
\[
\Gamma_{\beta,p}\min\{1,y_{n,d}\}
\le
C_{\beta,p}\min\{1,x_{n,d}\}.
\] % lean: degree_frontier

\item
Let \(R_{n,\ell_1}^{\star}(d,\beta,p,B)\) denote the coefficient-mass minimax risk in \cref{obj:def:two-class-minimax-risks}.  \Cref{obj:thm:bounded-outcome-degree-frontier} gives
\[
\alpha_{\beta,p}B^2\min\{1,y_{n,d}\}
\le
R_{n,\ell_1}^{\star}(d,\beta,p,B),
\qquad
R_{n,\ell_1}^{\star}(d,\beta,p,B)=R_n^\star(d,\beta,p,B).
\]
Combining this with the first rescaling yields
\[
c_{\beta,p}B^2\min\{1,x_{n,d}\}
\le
R_n^\star(d,\beta,p,B),
\]
which is the lower bound. % lean: degree_frontier

\item
The same application of \cref{obj:thm:bounded-outcome-degree-frontier} gives the projected-SNIPE upper bound at the energy scale,
\[
\sup_{(G,c)\in\mathcal M_{n,d,\beta}(B)}
\mathbb E_Z\!\left[
\left(\widehat\tau_n^{\mathrm{up}}-\tau_n\right)^2
\right]
\le
\Gamma_{\beta,p}B^2\min\{1,y_{n,d}\}.
\]
Using the second rescaling gives
\[
\sup_{(G,c)\in\mathcal M_{n,d,\beta}(B)}
\mathbb E_Z\!\left[
\left(\widehat\tau_n^{\mathrm{up}}-\tau_n\right)^2
\right]
\le
C_{\beta,p}B^2\min\{1,x_{n,d}\}.
\] % lean: degree_frontier

\item
It remains to place the minimax risk below the projected-SNIPE risk.  The model class is nonempty: the graph with no edges and the identically zero coefficient schedule belongs to \(\mathcal M_{n,d,\beta}(B)\), since its degree and coefficient mass are both zero and \(0\le B\).  Also, for every \((G,c)\in\mathcal M_{n,d,\beta}(B)\), every unit \(i\), and every assignment \(z\),
\[
\begin{aligned}
|Y_i(z)|
&=
\left|
\sum_{S\subseteq N_i}
c_{i,S}\prod_{j\in S}z_j
\right|  \\
&\le
\sum_{S\subseteq N_i}
|c_{i,S}|
\left|\prod_{j\in S}z_j\right|
\le
\sum_{S\subseteq N_i}|c_{i,S}|
\le B,
\end{aligned}
\]
because each \(z_j\in\{0,1\}\).  Hence, since \(n\ge1\),
\[
|\tau_n|
=
\left|
\frac1n\sum_{i\in V_n}
\{Y_i(\mathbf 1)-Y_i(\mathbf 0)\}
\right|
\le 2B.
\]
The projection defining \(\widehat\tau_n^{\mathrm{up}}\) gives
\[
|\widehat\tau_n^{\mathrm{up}}|\le B,
\]
and therefore
\[
\left(\widehat\tau_n^{\mathrm{up}}-\tau_n\right)^2
\le 9B^2
\]
pointwise.  Finite-space measurability is immediate, so \(\widehat\tau_n^{\mathrm{up}}\) is an admissible estimator for the coefficient-mass minimax problem. % lean: snipeClipped_admissible

\item
Since \(R_n^\star(d,\beta,p,B)\) is the infimum of worst risks over admissible estimators and the preceding step shows that \(\widehat\tau_n^{\mathrm{up}}\) is admissible,
\[
R_n^\star(d,\beta,p,B)
\le
\sup_{(G,c)\in\mathcal M_{n,d,\beta}(B)}
\mathbb E_Z\!\left[
\left(\widehat\tau_n^{\mathrm{up}}-\tau_n\right)^2
\right].
\] % lean: minimaxRisk_le_worstRisk_of_admissible

\item
For the unprojected SNIPE estimator, \cref{obj:thm:bounded-outcome-degree-frontier} gives
\[
\sup_{(G,c)\in\mathcal M_{n,d,\beta}(B)}
\mathbb E_Z\!\left[
\left(\widehat\tau_n^{\mathrm{SNIPE}}-\tau_n\right)^2
\right]
\le
B^2y_{n,d}.
\]
Using \(y_{n,d}\le\Gamma_{\beta,p}x_{n,d}\) and \(\Gamma_{\beta,p}\le C_{\beta,p}\),
\[
\sup_{(G,c)\in\mathcal M_{n,d,\beta}(B)}
\mathbb E_Z\!\left[
\left(\widehat\tau_n^{\mathrm{SNIPE}}-\tau_n\right)^2
\right]
\le
C_{\beta,p}B^2x_{n,d},
\]
which is the stated unprojected bound. % lean: degree_frontier

\item
For the unitwise energy statement, define, for \((G,c)\in\mathcal M_{n,d,\beta}(B)\) and \(i\in V_n\),
\[
A_i:=
\sum_{r=1}^{\min\{\beta,|N_i|\}}
\binom{|N_i|}{r}
\frac{\Delta_r(p)^2}{\{p(1-p)\}^r}.
\]
The local-energy conclusion of \cref{obj:thm:bounded-outcome-degree-frontier} gives
\[
A_i\le A_d,
\]
which is exactly
\[
\sum_{r=1}^{\min\{\beta,|N_i|\}}
\binom{|N_i|}{r}
\frac{\Delta_r(p)^2}{\{p(1-p)\}^r}
\le
A_d.
\] % lean: degree_frontier

\item
Fix \((G,c)\in\mathcal M_{n,d,\beta}(B)\), \(i\in V_n\), and \(r\in\mathbb N\) with \(1\le r\).  Put
\[
\mathcal S_{i,r}:=\{S\subseteq N_i:\ |S|=r\}.
\]
Since membership in \(\mathcal M_{n,d,\beta}(B)\) includes the bounded-degree condition with bound \(d\), \cref{obj:lem:overlap-count} gives
\[
\sum_{\ell\in V_n}
\binom{|N_i\cap N_\ell|}{r}
=
\sum_{S\in\mathcal S_{i,r}}
\#\{\ell\in V_n:S\subseteq N_\ell\},
\]
and
\[
\sum_{S\in\mathcal S_{i,r}}
\#\{\ell\in V_n:S\subseteq N_\ell\}
\le
d\binom{|N_i|}{r}
\le
d\binom{d}{r}.
\] % lean: overlap_count_le

\item
Finally assume \(0<B\) and \(1\le d\), and set
\[
m=\left\lfloor\frac nd\right\rfloor,
\qquad
\rho=\frac{md}{n},
\qquad
\delta=\delta(n,d,\beta,B,p).
\]
For every \(U:\{1,\ldots,m\}\to\mathbb R\) satisfying \(|U_b|\le B/2\), the block-mixture conclusion of \cref{obj:thm:bounded-outcome-degree-frontier} supplies models \(M_+\) and \(M_-\) in \(\mathcal M_{n,d,\beta}(B)\) with complete \(d\)-block graph and the corresponding sign-\(+1\) and sign-\(-1\) block schedules.  Their total treatment effects satisfy
\[
\tau_n(M_+)=\rho\delta,
\qquad
\tau_n(M_-)=-\rho\delta.
\]
The same conclusion gives the Hellinger bound
\[
H^2(\Pi_+,\Pi_-)
\le
\frac{4\pi^2m\delta^2}{B^2A_d},
\]
and the active-fraction identities
\[
\frac12\le\rho\le1,
\qquad
\frac{A_d}{m}
=
\frac{dA_d/n}{\rho}.
\]
These are precisely the remaining assertions. % lean: degree_frontier
\end{enumerate}
\end{proof}
